\begin{table}[htbp!]
\centering

% ---------- Table 3 ----------
\caption[TU traversal primitives.]{TU traversal primitives used to program fiber iteration. In each signature, \texttt{int} arguments are constants and \texttt{stream} arguments are values produced by a parent traversal. \texttt{Dense} implements a regular loop with static bounds, \texttt{Range} implements a loop with data-dependent bounds from streams, and \texttt{Index} implements a fixed-size dense lookup from a streamed base position. Overall, these primitives cover the loop-bound patterns needed for dense, compressed, and indirect traversal.}
\label{tab:tmu-traversal-primitives}
\resizebox{\textwidth}{!}{%
\begin{tabular}{@{}llll@{}}
\toprule
\textbf{Type} & \textbf{Description} & \textbf{Signature} & \textbf{Traversal}\\
\midrule
\texttt{Dense} & \makecell[l]{dense (or single)\\fiber scan} & \makecell[l]{\texttt{DnsFbrT(int beg, int end,}\\\texttt{ int stride=1)}} & \makecell[l]{for(i=beg; i<end;\\ \; i+=stride)}\\
\texttt{Range} & \makecell[l]{compressed fiber\\lookup and scan} & \makecell[l]{\texttt{RngFbrT(stream beg, stream end,}\\\texttt{ int offset=0, int stride=1)}} & \makecell[l]{for(i=beg.head()+offset;\\ \; i<end.head(); i+=stride)}\\
\texttt{Index} & \makecell[l]{dense fiber\\lookup and scan} & \makecell[l]{\texttt{IdxFbrT(stream beg, int size,}\\\texttt{ int offset=0, int stride=1)}} & \makecell[l]{for(i=beg.head()+offset;\\ \; i<beg.head()+size; i+=stride)}\\
\bottomrule
\end{tabular}
}

\vspace{1.2em}

% ---------- Table 4 ----------
\caption[TU data streams.]{Data streams produced or transformed by a TU. Memory streams load data using streamed indices, iteration and arithmetic streams generate addresses and loop positions, and forwarding or mask streams communicate values and predicates to the next layer. In this way, streams provide the dataflow interface between traversal, memory access, merging, and marshaling.}
\label{tab:tmu-data-streams}
\begin{tabular}{@{}llc@{}}
\toprule
\textbf{Type} & \textbf{Description} & \textbf{Function}\\
\midrule
\texttt{mem} & loads \texttt{p}'s elements from memory at index \texttt{x} & \texttt{p[x]}\\
\texttt{ite} & generates TU iteration indices $\in$[\texttt{beg},\texttt{end}) & \texttt{x:b->e}\\
\texttt{lin} & implements a linear transformation on \texttt{x} & \texttt{ax+b}\\
\texttt{map} & reads \texttt{x}-th element of a small map \texttt{a=\{v1,...,v16\}} & \texttt{a[x]}\\
\texttt{ldr} & computes the address of the \texttt{x}-th element in \texttt{p} & \texttt{\&p[x]}\\
\midrule
\texttt{fwd} &\multicolumn{2}{l}{forwards a leftward TU stream to the next layer}\\
\texttt{msk} &\multicolumn{2}{l}{outputs a layer predicate for merging and lockstep operations}\\
\bottomrule
\end{tabular}

\vspace{1.2em}

% ---------- Table 5 ----------
\caption[TMU inter-layer configurations.]{Inter-layer configurations used by Traversal Groups. \texttt{Single}, \texttt{BCast}, and \texttt{Keep} control how lanes are selected or broadcast, whereas \texttt{DisjMrg}, \texttt{ConjMrg}, and \texttt{LockStep} implement disjunctive merging, conjunctive merging, and lockstep co-iteration. In this way, these modes determine how lanes synchronize and how sparse fibers are merged across layers.}
\label{tab:tmu-layer-modes}
\begin{tabular}{@{}ll@{}}
\toprule
\textbf{Type} & \textbf{Description}\\
\midrule
\texttt{Single} & iterates a single lane\\
\texttt{BCast} & broadcasts a single lane to a parallel group\\
\texttt{Keep} & keeps one lane out of a parallel group\\
\midrule
\texttt{DisjMrg} & joins lanes of a layer\\
\texttt{ConjMrg} & intersects lanes of a layer\\
\texttt{LockStep} & co-iterates lanes of a layer\\
\bottomrule
\end{tabular}

\end{table}
